\section{The Milnor-Wood inequality}
Let $\Gamma$ be a uniform lattice in $\operatorname{SU}(1,n)$, $X$ be the compact quotient $\Gamma\backslash\mathbb H^n_{\mathbb C}$, and let $\rho$ be a representation of $\Gamma$ in a simple real algebraic Hermitian Lie group $G_{\mathbb R}$, whose associated symmetric space is denoted by $M$. In this section we use the material developed or recalled in Sections~2 and~3 to prove the Milnor-Wood inequality
\[
|\tau(\rho)|\leqslant r_M\operatorname{vol}(X).
\]
Before going into the proof, we remark that if the lattice $\Gamma$ has torsion, the measure defined by the Kähler form $\omega$ on $\mathbb H^n_{\mathbb C}$ descends to $X$ so that $\operatorname{vol}(X)$ is well-defined. There are several ways to define the Toledo invariant of the representation $\rho$ of $\Gamma$ in this case. Since this is how we will handle lattices with torsion, we observe that by Selberg's Lemma (see \cite[p.~327]{jep93Rat06} or \cite{jep93Sel60}) there exists a finite index torsion free subgroup $\Gamma'$ of $\Gamma$. We then define the Toledo invariant of $\rho:\Gamma\to G_{\mathbb R}$ by $\tau(\rho):=\tau(\rho')/|\Gamma/\Gamma'|$, where $\rho'=\rho\jepPubRestrict_{\Gamma'}$ is the restriction of $\rho$ to $\Gamma'$ and $\tau(\rho')$ is defined as in the introduction. One checks easily that this does not depend on the choice of $\Gamma'$. Moreover, if $X'=\Gamma'\backslash\mathbb H^n_{\mathbb C}$, then $\operatorname{vol}(X)=\operatorname{vol}(X')/|\Gamma/\Gamma'|$ so that proving the Milnor-Wood inequality for $\rho:\Gamma\to G_{\mathbb R}$ is equivalent to proving the Milnor-Wood inequality for $\rho':\Gamma'\to G_{\mathbb R}$.

In addition, the representation $\rho$ can always be deformed to a reductive representation with the same Toledo invariant, see e.g. \cite[Lem.~4.11]{jep93KM17}. Moreover, we may change the complex structure of $M$ to its opposite (i.e., replace $z\in\mathfrak t$ by $-z\in\mathfrak t$, or equivalently $\mathfrak m^{\pm}$ by $\mathfrak m^{\mp}$), thereby changing the sign of the Toledo invariant of $\rho$.

Therefore we may assume without loss of generality that
\begin{itemize}
\item the lattice $\Gamma$ is torsion free, so that $X=\Gamma\backslash\mathbb H^n_{\mathbb C}$ is a compact complex hyperbolic manifold of (complex) dimension $n$;
\item the representation $\rho$ is reductive, so that the results discussed in Section~3 are available;
\item the Toledo invariant $\tau(\rho)$ is positive, so we only have to prove $\tau(\rho)\leqslant r_M\operatorname{vol}(X)$.
\end{itemize}
(Observe that this implies that the harmonic map $f$ is not antiholomorphic.)

We also assume as in Section~2 that
\begin{itemize}\item the complexification $G$ of $G_{\mathbb R}$ is simply connected.\end{itemize}
See Remarks~5.14 and~5.15 where we explain why this is also not a loss of generality.
\subsection{Setup}
Consider the representation $\jepPubBbb{E}$ of $G$ defined in Section~2.2. As explained in Section~3.1, this gives rise to a flat harmonic $G_{\mathbb R}$-Higgs vector bundle $(\overline E,\overline\theta)$ over $X$, associated to the $G_{\mathbb R}$-Higgs principal bundle $(\overline P_K,\overline\theta)$. As a representation of $K$, $\jepPubBbb{E}=\jepPubOplus_{i=0}^{r_M}\jepPubBbb{E}_i$, where $r_M$ is the real rank of $G_{\mathbb R}$. This means that the holomorphic bundle $\overline E$ admits the holomorphic decomposition $\overline E=\jepPubOplus_{i=0}^{r_M}\overline E_i$, which is orthogonal for the harmonic metric. Moreover, the components $\overline\beta\in\overline P_K(\mathfrak m^-)\otimes\Omega_X^1$ and $\overline\gamma\in\overline P_K(\mathfrak m^+)\otimes\Omega_X^1$ of the Higgs field $\overline\theta\in\operatorname{Hom}(\overline E,\overline E)\otimes\Omega_X^1$ (see Remark~3.1) satisfy
\[
\overline\beta\in\jepPubOplus_{i=0}^{r_M-1}\operatorname{Hom}(\overline E_i,\overline E_{i+1})\otimes\Omega_X^1\quad\text{and}\quad\overline\gamma\in\jepPubOplus_{i=0}^{r_M-1}\operatorname{Hom}(\overline E_{i+1},\overline E_i)\otimes\Omega_X^1.
\]
We pull-back the harmonic Higgs bundle $(\overline E,\overline\theta)\to X$ to the projectivized tangent bundle $\mathbb PT_X$ of $X$ to obtain a harmonic Higgs bundle that we denote by $(E,\theta)\to\mathbb PT_X$. We restrict the Higgs field $\theta$ to the tangent space $L$ of the tautological foliation on $\mathbb PT_X$, so that its components $\beta$ and $\gamma$ satisfy
\[
\beta\in\jepPubOplus_{i=0}^{r_M-1}\operatorname{Hom}(E_i\otimes L,E_{i+1})\quad\text{and}\quad\gamma\in\jepPubOplus_{i=0}^{r_M-1}\operatorname{Hom}(E_{i+1}\otimes L,E_i).
\]
We also call $P_K$ the pull-back of the principal bundle $\overline P_K$ to $\mathbb PT_X$.
\begin{notation}
A tangent vector $\xi\neq0$ to $X$ at a point $x\in X$ defines a point $[\xi]$ in the projectivized tangent bundle $\mathbb PT_X$. It also defines an element in the fiber $\jepPubScript{O}_{\mathbb PT_X}(-1)_{[\xi]}=\{v\in T_{X,x}\mid v\in[\xi]\}$ of the tautological line bundle $\jepPubScript{O}_{\mathbb PT_X}(-1)$ at $[\xi]$, hence an element in the fiber $L_{[\xi]}$ of the tangent line bundle $L$ to the tautological foliation $\jepPubScript{T}$ at $[\xi]$, which will also be denoted by $\xi$.
\end{notation}
\begin{definition}
For $\xi\neq0$ a vector tangent to $X$ and $[\xi]$ the corresponding point in $\mathbb PT_X$, the \emph{rank} $\operatorname{rk}\beta_{[\xi]}$ of $\beta_{[\xi]}$ is the rank of the corresponding element $\beta(\xi)$ in $\mathfrak m^-$, as defined in Definition~2.22. By Proposition~2.24, it is also the largest value of $k$ such that $(\beta_{[\xi]})^k:E_0\otimes L^k\to E_k$ is not zero.
\begin{itemize}
\item The \emph{generic rank} $\operatorname{rk}\beta$ of $\beta$ is the maximum of the ranks of $\beta_{[\xi]}$ for $[\xi]\in\mathbb PT_X$.
\item The \emph{singular locus} of $\beta$ is the following subset of $\mathbb PT_X$:
\begin{align*}
\jepPubScript{S}(\beta)&:=\{[\xi]\in\mathbb PT_X\mid\operatorname{rk}\beta_{[\xi]}<\operatorname{rk}\beta\}\\
&=\{[\xi]\in\mathbb PT_X\mid(\beta_{[\xi]})^{\operatorname{rk}\beta}:E_0\otimes L^{\operatorname{rk}\beta}\to E_{\operatorname{rk}\beta}\text{ vanishes}\}.
\end{align*}
\item The \emph{regular locus} of $\beta$ is $\jepPubScript{R}(\beta):=\mathbb PT_X\backslash\jepPubScript{S}(\beta)$.
\item The \emph{singular locus} $\jepPubScript{S}(\overline\beta)$ of $\overline\beta$ is the projection to $X$ of $\jepPubScript{S}(\beta)$.
\end{itemize}
The \emph{regular locus} $\jepPubScript{R}(\overline\beta)$ of $\overline\beta$ is $X\backslash\jepPubScript{S}(\overline\beta)$ (it is not the projection to $X$ of $\jepPubScript{R}(\beta)$).
\end{definition}
Our assumption that $\tau(\rho)>0$ implies that $\overline\beta\neq0$ and $\beta\neq0$, therefore $1\leqslant\operatorname{rk}\beta\leqslant r_M$. Observe that while $\jepPubScript{S}(\beta)$ is an analytic subset of $\mathbb PT_X$ of codimension at least 1, its projection $\jepPubScript{S}(\overline\beta)$, although it is an analytic subset of $X$ (because $\pi:\mathbb PT_X\to X$ is a proper map), might be equal to the whole base $X$. A very important consequence of the weak polystability along the leaves (Proposition~3.2~(2)), is that this does not happen when there is equality in the refined Milnor-Wood inequality of Theorem~4.6, see Proposition~4.7. This will indeed be a crucial point when dealing with maximal representations.
\subsection{Rewording of the inequality}
Since the Hermitian symmetric space $M$ associated with $G_{\mathbb R}$ is a Kähler-Einstein manifold, the first Chern form $c_1(T_M)$ of its tangent bundle is a constant multiple of the $G_{\mathbb R}$-invariant Kähler form $\omega_M$: $c_1(T_M)=-\frac1{4\pi}c\,\omega_M$ for some positive constant $c$ ($c=c_1(\check M)$ is the first Chern number of the compact dual $\check M$ of $M$). On the other hand, the line bundle $\jepPubScript{L}$ associated with the $K$-representation $\jepPubBbb{E}_0$ is a generator of the Picard group of the compact dual $\check M$ of $M$ and it can be checked that the canonical bundle $K_{\check M}$ of $\check M$ is precisely given by $\jepPubScript{L}^{-c}$, see e.g. \cite[\S2]{jep93KM10}. Therefore the pull-back $f^\star\omega_M$ is $4\pi$ times the first Chern form of the line bundle $f^\star\jepPubScript{L}=\overline E_0$, so that the Toledo invariant of $\rho$ is
\begin{equation}\tag{4.1}
\tau(\rho)=4\pi\deg(\overline E_0)=4\pi\deg_{\jepPubScript{T}}(E_0),
\end{equation}
where the last equality follows from Proposition~3.3. Similarly, we get that
\begin{equation}\tag{4.2}
\deg(K_X)=\frac{n+1}{4\pi}\operatorname{vol}(X).
\end{equation}
Recall that $L$ is the tangential line bundle to the tautological foliation $\jepPubScript{T}$ on the projectivized tangent bundle $\mathbb PT_X$, and let $L^\star$ be its dual line bundle. One can compute as explained in \cite[\S4.3.1]{jep93KM17} that
\begin{equation}\tag{4.3}
\deg_{\jepPubScript{T}}(L^\star)=\frac1{2\pi}\operatorname{vol}(X).
\end{equation}
Therefore, the Milnor-Wood inequality can be rephrased as an inequality between the foliated degrees of the line bundles $E_0$ and $L^\star$ on $\mathbb PT_X$ and we need to prove:
\begin{equation}\tag{4.4}
\deg_{\jepPubScript{T}}(E_0)\leqslant\frac{r_M}{2}\deg_{\jepPubScript{T}}(L^\star),
\end{equation}
where $r_M$ is the rank of the symmetric space $M$.
\subsection{Leafwise Higgs subsheaves associated with the holomorphic component of the Higgs field}
We now define a subsheaf $\jepPubScript{V}$ of $\jepPubScript{E}:=\jepPubScript{O}(E)$ associated with $\beta$ (recall that $\beta\neq0$) in the same way we defined the submodule $\jepPubBbb{V}^r$ of $\jepPubBbb{E}$ associated with the nilpotent element $y\in\mathfrak m^-$ of rank $r$ in Definition~2.25 (for all $\xi\in L$, $\beta(\xi)\in P_K(\mathfrak m^-)$ is a nilpotent endomorphism of the bundle $E$). This subsheaf will be shown to be a leafwise Higgs subsheaf of the Higgs bundle $(E,\theta)$ on $\mathbb PT_X$. In Section~4.4 this will be used to prove the Milnor-Wood inequality.

More precisely, mimicking Definition~2.25 and for $-\operatorname{rk}\beta\leqslant k\leqslant\operatorname{rk}\beta$, we define the following subsheaves of $\jepPubScript{E}$:
\[
\jepPubScript{W}_k:=\sum_{\substack{\ell\geqslant0\\k+\ell+1\geqslant0}}\operatorname{Ker}\beta^{k+\ell+1}\cap\operatorname{Im}\beta^\ell,
\]
where in order to define $\operatorname{Ker}\beta^j$ we consider $\beta^j$ as a sheaf morphism from $\jepPubScript{E}$ to $\jepPubScript{E}\otimes(L^\star)^j$ and to define $\operatorname{Im}\beta^j$ we consider $\beta^j$ as a sheaf morphism from $\jepPubScript{E}\otimes L^j$ to $\jepPubScript{E}$.

For $k=0,1,\ldots,\operatorname{rk}\beta$, let $\jepPubScript{V}_k$ be the saturation in $\jepPubScript{E}_k:=\jepPubScript{O}(E_k)$ of the subsheaf $\jepPubScript{E}_k\cap\jepPubScript{W}_{\operatorname{rk}\beta-2k}$ and set
\[
\jepPubScript{V}=\jepPubOplus_{0\leqslant k\leqslant\operatorname{rk}\beta}\jepPubScript{V}_k.
\]
Since the sheaves $\jepPubScript{V}_k$ are saturated subsheaves of $\jepPubScript{O}(E_k)$, there exist a big open subset $\jepPubScript{U}$ of $\mathbb PT_X$ and subbundles $V_k$ of $E_k$ defined on $\jepPubScript{U}$ such that the restriction of the $\jepPubScript{V}_k$'s to $\jepPubScript{U}$ are the sheaves of sections of the $V_k$'s. On $\jepPubScript{U}$ we let $V$ be the subbundle $\oplus_{0\leqslant k\leqslant r}V_k$, so that $\jepPubScript{V}\jepPubRestrict_{\jepPubScript{U}}=\jepPubScript{O}_{\jepPubScript{U}}(V)$.

Observe that on the regular locus $\jepPubScript{R}(\beta)$ of $\beta$, the rank of $\beta^k$, as a vector bundle morphism from $E\otimes L^k$ to $E$, is constant for all $k$. Indeed, by Definition~2.22, the rank of $\beta(\xi)$ viewed as an element of $\mathfrak m^-$ determines its $K$-orbit. Hence on this open subset the formulas used above to define the subsheaves $\jepPubScript{W}_k$ of $\jepPubScript{E}$ in fact define subbundles $W_k$ of $E$ such that $\jepPubScript{W}_k\jepPubRestrict_{\jepPubScript{R}(\beta)}=\jepPubScript{O}_{\jepPubScript{R}(\beta)}(W_k)$. Therefore, on $\jepPubScript{R}(\beta)$, the subbundles $V_k$ such that $\jepPubScript{V}_k\jepPubRestrict_{\jepPubScript{R}(\beta)}=\jepPubScript{O}(V_k)$ are given by $V_k=E_k\cap W_{\operatorname{rk}\beta-2k}$. Thus, we may assume that $\jepPubScript{R}(\beta)$ is contained (and therefore dense) in $\jepPubScript{U}$.

We view an element $p$ of the $K$-principal bundle $P_K\overset{\pi_K}{\longrightarrow}\mathbb PT_X$ above $\xi\in\mathbb PT_X$ as an isomorphism between the fiber $E_\xi$ of $E=P_K(\jepPubBbb{E})$ and the model space $\jepPubBbb{E}$. The component $\beta$ of the Higgs field is a section of $P_K(\mathfrak m^-)\otimes L^\star\subset P_K(\operatorname{End}(\jepPubBbb{E}))\otimes L^\star$.

Let $y=y^{\operatorname{rk}\beta}$ and $Q_{\operatorname{rk}\beta}$ be defined as in Section~2.4.
\begin{lemma}
On the big open set $\jepPubScript{U}$, the subsheaf $\jepPubScript{V}$ defines a reduction $P_{K\cap Q_{\operatorname{rk}\beta}}$ of the structure group of $P_K$ to the subgroup $K\cap Q_{\operatorname{rk}\beta}\subset K$.
\end{lemma}
\begin{proof}
We begin by working on $\jepPubScript{R}(\beta)\subset\jepPubScript{U}$. Since for all $\xi\in\jepPubScript{R}(\beta)$ and all $\eta\in L_\xi$, $\eta\neq0$, we have that $\beta_\xi(\eta)$ has rank $\operatorname{rk}\beta$, there exists $p\in(P_K)_\xi$ such that $p\circ\beta_\xi(\eta)\circ p^{-1}=y\in\mathfrak m^-\subset\operatorname{End}(\jepPubBbb{E})$, so that $p(V_\xi)=\jepPubBbb{V}^{\operatorname{rk}\beta}$. Therefore, on $\jepPubScript{R}(\beta)$, by Proposition~2.33~(4), the subbundle $V$ of $E$ defines a (holomorphic) reduction $P_{K\cap Q_{\operatorname{rk}\beta}}$ of the structure group of $P_K$ to the subgroup $K\cap Q_{\operatorname{rk}\beta}$ of $K$ ($Q_{\operatorname{rk}\beta}$ is the normalizer in $G$ of the parabolic subalgebra $\mathfrak q_{\operatorname{rk}\beta}$, see Definition~2.30). Explicitly $P_{K\cap Q_{\operatorname{rk}\beta}}=\{p\in P_K\mid p(V_{\pi_K(p)})=\jepPubBbb{V}^{\operatorname{rk}\beta}\}$.

We now work on $\jepPubScript{U}$. Enlarge the structure group of $P_K$ to $\operatorname{GL}(\jepPubBbb{E})$. The subbundle $V=\jepPubOplus_{k=0}^{\operatorname{rk}\beta}V_k$ of $E$ defines a reduction $P_S$ of the structure group of $P_{\operatorname{GL}(\jepPubBbb{E})}\overset{\pi_{\operatorname{GL}(\jepPubBbb{E})}}{\longrightarrow}\mathbb PT_X$ to the stabilizer $S$ of $\jepPubBbb{V}^{\operatorname{rk}\beta}$ in $\operatorname{GL}(\jepPubBbb{E})$ by setting $P_S=\{p\in P_{\operatorname{GL}(\jepPubBbb{E})}\mid p(V_{\pi_{\operatorname{GL}(\jepPubBbb{E})}(p)})=\jepPubBbb{V}^{\operatorname{rk}\beta}\}$.

Let $B\subset\jepPubScript{U}$ be an open ball on which $P_K$ is trivial. Then the reductions $P_{K\cap Q_{\operatorname{rk}\beta}}$ of $P_K$ on $B\cap\jepPubScript{R}(\beta)$ and $P_S$ of $P_{\operatorname{GL}(\jepPubBbb{E})}$ on $B$ are respectively given by holomorphic maps $\sigma:B\cap\jepPubScript{R}(\beta)\to K/(K\cap Q_{\operatorname{rk}\beta})$ and $s:B\to\operatorname{GL}(\jepPubBbb{E})/S$. Moreover, if $\iota$ denotes the natural map $K/(K\cap Q_{\operatorname{rk}\beta})\to\operatorname{GL}(\jepPubBbb{E})/S$, which is injective, then we have $s=\iota\circ\sigma$ on $B\cap\jepPubScript{R}(\beta)$. Since $K/(K\cap Q_{\operatorname{rk}\beta})$ is compact, its image by $\iota$ is closed in $\operatorname{GL}(\jepPubBbb{E})/S$. Therefore, since $B\cap\jepPubScript{R}(\beta)$ is dense in $B$, $s$ maps $B$ to $\iota(K/(K\cap Q_{\operatorname{rk}\beta}))$. This means that the reduction $P_{K\cap Q_{\operatorname{rk}\beta}}$ initially defined on $\jepPubScript{R}(\beta)$ extends to $\jepPubScript{U}$.
\end{proof}
We deduce the following result.
\begin{proposition}
The subsheaf $\jepPubScript{V}$ is a leafwise Higgs subsheaf of the Higgs bundle $(E,\theta)$ on $\mathbb PT_X$.
\end{proposition}
\begin{proof}
By Proposition~2.29, we know that $y$ and $\mathfrak m^+$ stabilize $\jepPubBbb{V}^{\operatorname{rk}\beta}$. Therefore, on $\jepPubScript{R}(\beta)$, the two components $\beta$ and $\gamma$ of the Higgs field stabilize the subsheaf $\jepPubScript{V}$ since it is the sheaf of sections of the subbundle $V=P_{K\cap Q_{\operatorname{rk}\beta}}(\jepPubBbb{V}^{\operatorname{rk}\beta})$ of $E=P_{K\cap Q_{\operatorname{rk}\beta}}(\jepPubBbb{E})$. By continuity, this still holds on $\jepPubScript{U}$ since on this big open set $\jepPubScript{V}$ is also the sheaf of section of $V=P_{K\cap Q_{\operatorname{rk}\beta}}(\jepPubBbb{V}^{\operatorname{rk}\beta})$. Now, $\jepPubScript{V}$ is a saturated, hence normal, subsheaf of $\jepPubScript{O}(E)$ by definition. Hence the restriction map $\jepPubScript{V}(\mathbb PT_X)\to\jepPubScript{V}(\jepPubScript{U})$ is an isomorphism since $\jepPubScript{U}$ is big. Therefore $\jepPubScript{V}$ is indeed a leafwise Higgs subsheaf of $(E,\theta)$ on $\mathbb PT_X$.
\end{proof}
\subsection{Proof of the Milnor-Wood inequality}
Recall that the slope of an $H$-module was introduced in Section~2.5 and is an element in $X(H)\otimes\mathbb Q$. To relate it to the slope of the associated vector bundle, we use the following fact which follows immediately from the definitions.
\begin{fact}
Let $H$ be a complex reductive group and let $P_H\to Z$ a holomorphic $H$-principal bundle over a complex manifold $Z$. Let $\jepPubBbb{V}_i,i\in\{1,2\}$, be $H$-modules and let $V_i=P_H(\jepPubBbb{V}_i)$ be the associated vector bundles. Assume that $\mu_H(\jepPubBbb{V}_1)=\mu_H(\jepPubBbb{V}_2)$. Then
\[
\frac{\det V_1}{\operatorname{rk}V_1}=\frac{\det V_2}{\operatorname{rk}V_2},
\]
as elements of $\operatorname{Pic}(Z)\otimes_{\mathbb Z}\mathbb Q$. In other words, the line bundles $(\det V_1)^{\otimes\operatorname{rk}V_2}$ and $(\det V_2)^{\otimes\operatorname{rk}V_1}$ are isomorphic.
\end{fact}
Therefore, if one is given a reasonable way to compute slopes of vector bundles on $Z$, e.g. the usual one if $Z$ is $X$ with its Kähler form $\omega$, or the foliated one if $Z$ is $\mathbb PT_X$ with its tautological foliation $\jepPubScript{T}$ and transverse volume form $\Omega_{\jepPubScript{G}}$, then using the notation of Fact~4.5, $\mu_H(\jepPubBbb{V}_1)=\mu_H(\jepPubBbb{V}_2)$ implies that the slopes of $V_1$ and $V_2$ are equal.

This observation, together with the computation of the slopes of the $H_{\operatorname{rk}\beta}$-submodules $\jepPubBbb{V}^{\operatorname{rk}\beta}_k$ of $\jepPubBbb{E}$ and the construction of the leafwise Higgs subsheaf $\jepPubScript{V}$ of $(E,\theta)$, gives the Milnor Wood inequality.
\begin{theorem}
We have the inequalities
\[
\deg_{\jepPubScript{T}}(E_0)+\frac{\operatorname{rk}\beta}{2}\deg_{\jepPubScript{T}}(L)\leqslant\mu_{\jepPubScript{T}}(\jepPubScript{V})\leqslant0.
\]
Therefore the Milnor-Wood inequality $\deg_{\jepPubScript{T}}(E_0)\leqslant(r_M/2)\deg_{\jepPubScript{T}}(L^\star)$ holds.
\end{theorem}
\begin{proof}
Let $n_k:=\dim\jepPubBbb{V}^{\operatorname{rk}\beta}_k$. First, recall that Proposition~2.33~(2) states that the unipotent radical of $Q_{\operatorname{rk}\beta}$ acts trivially on $\jepPubBbb{V}^{\operatorname{rk}\beta}$. Hence so does the unipotent radical of $K\cap Q_{\operatorname{rk}\beta}$. Thus, in fact, $\jepPubBbb{V}^{\operatorname{rk}\beta}$ is a $(K\cap Q_{\operatorname{rk}\beta})/R_u(K\cap Q_{\operatorname{rk}\beta})$-module, and $(K\cap Q_{\operatorname{rk}\beta})/R_u(K\cap Q_{\operatorname{rk}\beta})\simeq H_{\operatorname{rk}\beta}$ is reductive.

Therefore, by Lemma~4.3, on the big open set $\jepPubScript{U}$, the vector bundles $V_k$ such that $\jepPubScript{V}_k=\jepPubScript{O}(V_k)$ are given by $V_k=P_{K\cap Q_{\operatorname{rk}\beta}}(\jepPubBbb{V}^{\operatorname{rk}\beta}_k)$. By Proposition~2.45 and Fact~4.5 applied to $\jepPubBbb{V}^{\operatorname{rk}\beta}_k$ and $\jepPubBbb{V}'_k:=\jepPubBbb{V}^{\operatorname{rk}\beta}_0\otimes(\jepPubBbb{V}^{\operatorname{rk}\beta}_1\otimes\jepPubBbb{V}^{\operatorname{rk}\beta}_0{}^\star)^{\otimes k}$, we have on $\jepPubScript{U}$
\[
(\det V_k)^{\operatorname{rk}V'_k}\simeq(\det V'_k)^{\operatorname{rk}V_k},
\]
where $V'_k=P_{K\cap Q_{\operatorname{rk}\beta}}(\jepPubBbb{V}'_k)$. Letting $\jepPubScript{V}'_k:=\jepPubScript{V}_0\otimes(\jepPubScript{V}_1\otimes\jepPubScript{V}_0{}^\star)^{\otimes k}$, we deduce that the line bundles $(\det\jepPubScript{V}_k)^{\operatorname{rk}\jepPubScript{V}'_k}$ and $(\det\jepPubScript{V}'_k)^{\operatorname{rk}\jepPubScript{V}_k}$ on $\mathbb PT_X$ are isomorphic on $\jepPubScript{U}$, hence on $\mathbb PT_X$, because $\jepPubScript{U}$ is a big open set of $\mathbb PT_X$. Therefore
\[
\mu_{\jepPubScript{T}}(\jepPubScript{V}_k)=\mu_{\jepPubScript{T}}(\jepPubScript{V}'_k)=\deg_{\jepPubScript{T}}(\jepPubScript{V}_0)+k\,\mu_{\jepPubScript{T}}(\jepPubScript{V}^\star_0\otimes\jepPubScript{V}_1).
\]
Since $\beta^{\operatorname{rk}\beta}:\jepPubScript{V}_0\otimes L^{\operatorname{rk}\beta}\to\jepPubScript{V}_{\operatorname{rk}\beta}$ is not zero and $n_0=n_{\operatorname{rk}\beta}=1$, we also have $\mu_{\jepPubScript{T}}(\jepPubScript{V}_{\operatorname{rk}\beta})\geqslant\mu_{\jepPubScript{T}}(\jepPubScript{V}_0)+\operatorname{rk}\beta\,\mu_{\jepPubScript{T}}(L)$ so that $\mu_{\jepPubScript{T}}(\jepPubScript{V}^\star_0\otimes\jepPubScript{V}_1)\geqslant\deg_{\jepPubScript{T}}(L)$.

Finally, remembering that $n_k=n_{\operatorname{rk}\beta-k}$ by Proposition~2.35 and that $\jepPubScript{V}_0=E_0$, we get
\begin{align*}
2\deg_{\jepPubScript{T}}(\jepPubScript{V})&=\sum_{k=0}^{\operatorname{rk}\beta}\deg_{\jepPubScript{T}}(\jepPubScript{V}_k)+\sum_{k=0}^{\operatorname{rk}\beta}\deg_{\jepPubScript{T}}(\jepPubScript{V}_{\operatorname{rk}\beta-k})\\
&=\sum_{k=0}^{\operatorname{rk}\beta}(n_k\mu_{\jepPubScript{T}}(\jepPubScript{V}_k)+n_{\operatorname{rk}\beta-k}\mu_{\jepPubScript{T}}(\jepPubScript{V}_{\operatorname{rk}\beta-k}))\\
&\geqslant\sum_{k=0}^{\operatorname{rk}\beta}n_k(\deg_{\jepPubScript{T}}(\jepPubScript{V}_0)+k\deg_{\jepPubScript{T}}(L)+\deg_{\jepPubScript{T}}(\jepPubScript{V}_0)+(\operatorname{rk}\beta-k)\deg_{\jepPubScript{T}}(L))\\
&=(\dim\jepPubBbb{V}^{\operatorname{rk}\beta})(2\deg_{\jepPubScript{T}}(E_0)+\operatorname{rk}\beta\deg_{\jepPubScript{T}}(L)).
\end{align*}
The inequality $\mu_{\jepPubScript{T}}(\jepPubScript{V})\leqslant0$ follows from Propositions~4.4 and~3.2, and the conclusion from the fact that $\operatorname{rk}\beta\leqslant r_M$.
\end{proof}
In case of equality in Theorem~4.6, we have
\begin{proposition}
Assume that $\deg_{\jepPubScript{T}}(E_0)+\frac{\operatorname{rk}\beta}{2}\deg_{\jepPubScript{T}}(L)=0$. Then
\begin{enumerate}[label=(\arabic*)]
\item on the regular locus $\jepPubScript{R}(\beta)=\mathbb PT_X\backslash\jepPubScript{S}(\beta)$ of $\beta$, the orthogonal complement $V^\perp=(\jepPubOplus V_i)^\perp$ of the subbundle $V=\jepPubOplus V_i$ of $E$ w.r.t. the harmonic metric is stable under the Higgs field $\theta:E\otimes L\to E$;
\item the regular locus $\jepPubScript{R}(\overline\beta)\subset X$ of $\overline\beta$ is (open and) dense in $X$.
\end{enumerate}
\end{proposition}
\begin{proof}
The first point follows from the discussion after the definition of the subsheaves $\jepPubScript{V}_k$ and the polystability property~(2) in Proposition~3.2, since our hypothesis implies that $\deg_{\jepPubScript{T}}\jepPubScript{V}=0$ by Theorem~4.6.

Proposition~3.2~(2) also implies that the singular locus $\jepPubScript{S}(\beta)$ of $\beta$, which is a closed proper subset of $\mathbb PT_X$, is saturated under the tautological foliation $\jepPubScript{T}$, see the proof of \cite[Lem.~4.5]{jep93KM17}. Since the leaves of the tautological foliation on $\mathbb PT_X$ are projections of $\operatorname{SU}(1,1)$-homogeneous subsets of $\Gamma\backslash\operatorname{SU}(1,n)$, M.~Ratner's results on unipotent flows and the fact that $\mathbb H^n_{\mathbb C}=\operatorname{SU}(1,n)/\operatorname U(n)$ is a rank~$1$ symmetric space, then imply by \cite[Prop.~3.6]{jep93KM17} that the singular locus $\jepPubScript{S}(\overline\beta):=\pi(\jepPubScript{S}(\beta))$ of $\overline\beta$ is a proper analytic subset of $X=\Gamma\backslash\mathbb H^n_{\mathbb C}$, hence the second point of the proposition.
\end{proof}

